\documentstyle[%


righttag%


,11pt%


,amssymb%


,russian%               remove in Eng.


,izvru%                 Set izven% in Eng.


%,izvru-ks             for brief communications


]{amsart}





%





\newcommand{\diag}{\operatorname{diag}}


\newcommand{\const}{\operatorname{const}}


\newcommand{\re}{\operatorname{Re}}


\newcommand{\tts}[1]{}





\theoremstyle{plain}


\theorembodyfont{\it}


\newtheorem{thmnonum}{\hskip\parindent Теорема}


\renewcommand{\thethmnonum}{}





%\hoffset=-15mm%                  For Eng.


\setcounter{page}{36}


\god{2004}


\nomer{2~(501)} %                        Remove in Eng.


%% remove ^^^^^^


%\nomer{Vol.\,48, No.\,2}%               Set in Eng.


%\str{\pageref{begin}--\pageref{end}}%  Set in Eng.


\udk{517.968}





\author{\it Т.О.\,Колбинева, З.Б.\,Цалюк}


% NB:   ^^ remove in Eng.


\title{Асимптотическое поведение решений одного класса интегральных уравнений}





\date{}


%\pagestyle{myheadings}%         Set in Eng.


\pagestyle{plain} %              Remove in Eng.


\begin{document}


%\thispagestyle{plain}%          Set in Eng.


\maketitle


\label{begin}








Рассмотрим уравнение


\begin{equation}


\label{1}


x (t)=\int_0^t \sum\limits_{j=1}^{n} K_j(t-s) a_j(s)


x(s) ds+f(t),


\end{equation}


где $K_j$, $a_j$, $f$ --- непрерывные на $[0,\infty)$ комплекснозначные 


функции.





Если $a_j(t)=\const$ и $K_j \in L_1[0,\infty)$, то асимптотика


решений уравнения~(\ref{1}) хорошо изучена (см., напр., [1]--[4]). Для 


периодических функций $a_j(t)$ асимптотика решений изучалась в [5], [6]. В данной


статье рассмотрен случай, когда $a_j(t)$ имеют асимптотическое разложение 


по степенной шкале в окрестности~$+\infty$:


\begin{equation}


\label{2}


a_j(t)\sim a_{j0}+\frac{a_{j1} }{{t+1}}+\dotsb+\frac{{a_{jk}}}{{(t+1)^k }}+


\dotsb,\quad t \to \infty.


\end{equation}


(Использование степеней $(t+1)^k$ вместо $t^k$ вызвано лишь техническими


удобствами.)





Предположим, что при всех достаточно малых $\delta>0$ и некотором 


$\mu$


\begin{equation}


\label{3}


e^{-\delta t}K_j(t) \in L_1[0,\infty),\ \ K_j(t)=O (e^{\mu t}),


\ \ f(t)=O (e^{\mu t}).


\end{equation}


Обозначим через $\lambda_l$ расположенные в полуплоскости $\re z\ge 0$ 


корни уравнения


\begin{equation}


\label{4}


\sum\limits_{j=1}^{n} a_{j0} \wh K_j(z)=1,


\end{equation}


где $\wh K_j (z)=\int\limits_0^\infty e^{-zt}K_j (t)dt$ -- 


преобразование Лапласа, и пусть $\lambda_1, \dotsb, \lambda_r$ --- 


все корни уравнения~(\ref{4}), для которых


$\re \lambda_k=\lambda=\max\limits_l \re {\lambda_l }$.





Обозначим через $A_k=\Big\{ u(t) \in C[0, \infty): u(t)=


\sum\limits_{j=0}^k \frac {u_j} {(t+1)^j}+o (\frac 1 {(t+1)^k}


) \Big\}$ и $A_{\infty}=\bigcap\limits_{k=0}^{\infty}A_k$. Таким образом,


условие~(\ref{2}) означает, что $a_j \in A_{\infty}$.





Основной результат представляет





\begin{thmnonum}


Пусть выполнены условия~$(\ref{2})$ и~$(\ref{3})$ с 


$\mu <


\lambda$ и все корни $\lambda_1, \dotsc, \lambda_r$ простые.


Тогда решение $x(t)$ уравнения~$(\ref{1})$ представимо в виде


$$


x(t)=\sum\limits_{j=1}^r e^{\lambda_j t} (t+1)^{\alpha_j} b_j(t),


$$


где $\alpha_j$ -- некоторые комплексные числа, а $b_j \in A_{\infty}$.


\end{thmnonum}





Заметим, что при $r>1$ функции $b_j(t)$ определяются неоднозначно.


Например, функции $b_1$ и $b_2$ можно заменить на


$b_1 (t)+(t+1)^{-\alpha_1}e^{- (c+\lambda_1)t}$


и $b_2 (t)- (t+1)^{-\alpha_2}e^{- (c+\lambda_2)t}$,


$c>0$.


Однако можно показать, что при каждом $j$ коэффициенты асимптотического 


разложения различных $b_j(t)$ одинаковы.





Докажем сначала три леммы.





\begin{lem}


Уравнение~$(\ref{1})$ эквивалентно некоторому уравнению вида


\begin{equation}


\label{5}


x(t)=\int_0^t H(t,s)x(s)ds+\sum\limits_{j=1}^r \int_0^t G_j(t,s)


\int_0^s H(s,\tau)x(\tau)d\tau ds+


\sum\limits_{j=1}^r e^{\lambda_j t}


(t+1)^{\alpha_j} y_j(t),


\end{equation}


где


$H(t,s)=\sum\limits_{l=1}^n h_l(t-s) \frac {\gamma_l(s)}{s+1}$,


$G_j(t, s)=e^{\lambda_j(t-s)} (\frac {t+1} {s+1})^{\alpha_j}


u_j(t)\frac {v_j(s)}{s+1}$,


$\alpha_j$ --- некоторые комплексные числа, $\gamma_l,


u_j, v_j, y_j \in A_{\infty}$, $h_l(t)=o (e^{(\lambda-


\varepsilon)t})$.


\end{lem}








{\bf Доказательство.} Положим 


$K_{\delta}(t)=e^{-\delta t} \sum\limits_{j=1}^na_{j0}K_{j}(t)$, 


где $\delta \ge 0$. Пусть $\delta>0$, достаточно мало и таково, что


уравнение~(\ref{4}) не имеет корней на прямой $\re z=\delta$. Тогда


уравнение~(\ref{4}) имеет в полуплоскости $\re z \ge \delta$ конечное число


корней $\lambda_1, \dotsc, \lambda_m$ целой кратности. Очевидно, уравнение


$\wh K_{\delta}(z)=1$ имеет в полуплоскости $\re z \ge 0$ корни $\lambda_1-


\delta, \dotsc, \lambda_m-\delta$ и потому (см., напр., [1]) резольвента 


$R_{\delta}(t)$ ядра $K_{\delta }(t)$ представима в виде


$$


R_{\delta}(t)=R_0(t)+\sum\limits_{j=1}^r c_j e^{(\lambda_j-\delta)t}+


e^{(\lambda-\delta-\varepsilon_1)t} \phi (t)+


R_0(t) * \bigg[\sum\limits_{j=1}^r c_j


e^{(\lambda_j-\delta)t}+e^{(\lambda-\delta-\varepsilon_1)t} \phi (t)\bigg],


$$


где $h (t)* g(t)=\int\limits_0^t h(t-s)g(s)ds$, $R_0 \in L_1[0, \infty)$,


$0<\varepsilon_1<\lambda-\max\limits_{j>r} \re \lambda_j $, а


$\phi(t)=o(1)$.





Так как


$R_0(t) * e^{(\lambda_j-\delta)t}=c e^{(\lambda_j-\delta)t}+o(1)$,


а


$$


R_0 (t)* e^{(\lambda-\delta-\varepsilon_1)t} \phi (t)=e^{(\lambda-\delta-


\varepsilon_1)t} \int_0^t e^{-(\lambda-\delta-\varepsilon_1)(t-s)} 


R_0(t-s)\phi(s)ds


=e^{(\lambda-\delta-\varepsilon_1)t} o(1),


$$


то


$R_{\delta}(t)=\sum\limits_{j=1}^r c_j e^{(\lambda_j-\delta)t}+o (


e^{(\lambda-\delta-\varepsilon_1)t})+R_0(t)$.


Резольвента $R(t)$ ядра $K_{0}(t)=\sum\limits_{j=1}^n a_{j0}K_{j}(t)$ 


связана с $R_{\delta}(t)$ равенством $R(t)=e^{\delta t}R_{\delta}(t)$. 


Следовательно,


\begin{equation}


\label{6}


R(t)=\sum\limits_{j=1}^r c_j e^{\lambda_j t}+o (e^{(\lambda-


\varepsilon_1)t}


)+e^{\delta t}R_0(t).


\end{equation}


Записывая уравнение~(\ref{1}) в виде


$$


x(t)-K_0(t)*x(t)=\sum\limits_{j=1}^n K_j(t)* [ (a_j(t)-


a_{j0})x(t) ]+f(t)


$$


и используя резольвенту $R(t)$ ядра $K_0$, найдем


\begin{equation}


\label{7}


x(t)=\sum\limits_{j=1}^n [


K_j(t)+R(t)*K_j(t) ]* [ (a_j(t)-a_{j0})x(t) ]


+f(t)+R(t)*f(t).


\end{equation}





Заметим, что если $\delta<\mu$, $\mu+\varepsilon_1<\lambda$,


$\varepsilon<\frac{\lambda-\mu}2$ и $0<\varepsilon<\varepsilon_1$, то


$e^{\delta t}R_0(t)*e^{\mu t} O(1)=O (e^{\mu t})=o (e^{(\lambda-


\varepsilon) t})$,


\begin{gather*}


e^{\lambda_j t }*e^{\mu t} O (1)=e^{\lambda_j t }


\int_0^\infty e^{- (\lambda_j-\mu)s } O (1)ds-


e^{\lambda_j t } \int_t^\infty e^{-\frac{\lambda-\mu}2 s}


e^{-\frac{\lambda-\mu}2 s} O (1)ds=\\


%


=c e^{\lambda_j t }+e^{ (\lambda-\frac{\lambda-\mu}2) t}


o (1)=c e^{\lambda_j t}+o (e^{(\lambda-\varepsilon) t})


\end{gather*}


и


$e^{(\lambda-\varepsilon_1) t} o (1)* e^{\mu t} O (1)=


e^{(\lambda-\varepsilon_1) t}O (1)=


o (e^{(\lambda-\varepsilon)t})$.


Отсюда из~(\ref{3}),~(\ref{6}) и~(\ref{7}) получим


$$


x (t)=\sum\limits_{j=1}^r


e^{\lambda_j t }*\frac{\beta_j (t)}{t+1}x (t)+


\sum\limits_{l=1}^n h_l (t)* \frac{\gamma_l (t)}{t+1}


x (t)+\sum\limits_{j=1}^r k_j e^{\lambda_j t }+h_0 (t),


$$


где $\beta_j, \gamma_l \in A_{\infty}$, $h_l (t)=o (e^{(\lambda-


\varepsilon) t})$, $l=0,\dotsc,n$, $j=1,\dotsc,r$. Положим


\begin{gather*}


\psi_j (t)=k_j e^{\lambda_j t }+e^{\lambda_j t }*


\frac{\beta_j (t)}{t+1}x (t), \\


%


p (t)=h_0 (t)+\sum\limits_{l=1}^n h_l (t)*


\frac{\gamma_l (t)}{t+1} x (t)=h_0 (t)+


\int_0^t H (t,s) x (s) ds, \\


%


\psi (t)=


(\psi_1 (t),\dotsc,\psi_r (t))^T,


\ \ k=(k_1,\dotsc,k_r)^T,


\ \ \beta (t)=(\beta_1 (t),\dotsc,\beta_r (t))^T,


\end{gather*}


($a^T$ --- вектор-столбец, транспонированный $a$),


$\Lambda=\diag (\lambda_j)$ и $A (t)=(A_{jl} (t))$ с


$A_{jl} (t)=\beta_j (t)$. Тогда


$$


\psi(t)=e^{\Lambda t}k+\int_0^t e^{\Lambda (t-s)} \frac {A(s)}


{s+1}\psi(s)ds+\int_0^t e^{\Lambda (t-s)} \frac {\beta(s)} {s+1}


p(s)ds,


$$


откуда


\begin{equation}


\label{8}


\psi' (t)=\bigg(\Lambda+\frac {A(t)} {t+1} \bigg) \psi(t)+\frac


{\beta(t)} {t+1} p(t).


\end{equation}


По формуле Коши


$$


\psi(t)=Y(t)Y^{-1}(0) \psi(0)+\int_{0}^t


Y(t)Y^{-1}(s) \frac {\beta(s)} {s+1} p(s) ds,


$$


где $Y (t)$ --- фундаментальная матрица системы~(\ref{8}). В силу


известной теоремы (см., напр., [7], гл.\,V) в качестве $Y (t)$ 


можно выбрать матрицу вида


$Y (t)=D (t)\diag (e^{\lambda_j t}(t+1)^{\alpha_j})$


с невырожденной матрицей $D (t)\in A_{\infty}$, причем 


$\det D (\infty)=\det\lim\limits_{t\to\infty}


D (t) \neq 0$. Тогда $D^{-1} (t) \in A_{\infty}$ и


$v (t)=D^{-1} (t)\beta (t) \in A_{\infty}$.


Поэтому


$$


Y(t)Y^{-1}(s)\beta (s)=


D(t)\diag \bigg(e^{\lambda_j (t-s)} \bigg(\frac {t+1} {s+1} \bigg)^


{\alpha_j} \bigg)v (s)=


\begin{pmatrix}


\sum\limits_{j=1}^r e^{\lambda_j (t-s)} (\frac {t+1} {s+1})^{


\alpha_j }d_{1j} (t)v_j (s)\\


\vdots\\


\sum\limits_{j=1}^r e^{\lambda_j (t-s)} (\frac {t+1} {s+1})^{


\alpha_j }d_{rj} (t)v_j (s)


\end{pmatrix}


,


$$


где $ d_{kj}, v_j \in A_{\infty}$. Подставляя этот вектор в формулу Коши и


используя тот факт, что


\begin{gather*}


\int_0^t G_j (t,s) h_0 (s)ds=


e^{\lambda_j t} (t+1)^ {\alpha_j} u_j (t)


\bigg[\int_0^\infty e^{-\varepsilon s} o ((s+1)^{-


\alpha_j-1})ds - \\


%


-\int_t^\infty e^{-\varepsilon s} o ((s+1)^{-


\alpha_j-1})ds \bigg]=e^{\lambda_j t} (t+1)^ {\alpha_j}


u_j (t) [c_j+o (e^{-\frac \varepsilon 2 t}) ],


\end{gather*}


найдем (при $u_j (t)=\sum\limits_{k=1}^r d_{kj} (t)$)


\begin{gather*}


x (t)=\sum\limits_{k=1}^r \psi_k (t)+p (t)=


\sum\limits_{j=1}^r e^{\lambda_j t} (t+1)^ {\alpha_j}


y_j (t)+\\


%


+\int_0^t H (t,s)x (s)ds+


\sum\limits_{j=1}^r \int_0^t G_j (t,s)


\int_0^s H (s,\tau)x (\tau)d\tau ds.\quad\square


\end{gather*}





\begin{lem}


Пусть $h(t)=O (e^{-\varepsilon t})$,


$\varepsilon>0$, $\alpha$ -- комплексное число.


Тогда


\begin{gather*}


h(t)\ast (t+1)^{\alpha}=(t+1)^{\alpha}v (t),\ \ v \in A_\infty,\\


%


h(t)\ast O ((t+1)^ {\alpha})=O ((t+1)^ {\alpha}),


\ \ h(t)\ast o ((t+1)^ {\alpha})=o ((t+1)^ {\alpha}).


\end{gather*}


\end{lem}





\begin{pf}


Так как функция


$e^{\frac{-\varepsilon}2 (t-s)} (\frac{s+1}{t+1})^{\alpha}$


ограничена при $0\le s \le t < \infty$, то


$$


(t+1)^{-\alpha} [h(t)\ast o ((t+1)^ {\alpha}


) ]=O (1)\int_0^t e^{\frac{-\varepsilon}2 (t-


s)} o (1) ds=o (1),


$$


что доказывает последнее утверждение леммы. Аналогично доказывается и второе


утверждение. Наконец, интегрируя достаточное число раз по частям, получим


\begin{gather*}


h(t)\ast (t+1)^{\alpha}=\int_0^t (t-s+1)^{\alpha}


h (s) ds=\sum\limits_{j=0}^{k} c_j (t+1)^{\alpha-j}+\\


%


+\sum\limits_{j=0}^{k}b_j\int_t^\infty (t-\tau)^j


h (\tau) d\tau+c\int_{0}^t (t-s+1)^ {\alpha-k-1}


\int_s^\infty (\tau-s)^k h (\tau) d\tau ds.


\end{gather*}


Так как $\int\limits_t^\infty (t-


\tau)^{j}h (\tau)d\tau=\int\limits_0^\infty (-\theta)^{j} h (t+


\theta)d\theta=O (e^{{-\varepsilon}{t}})$,


то второе слагаемое в $ h(t)\ast (t+1)^{\alpha}$ есть


$O (e^{-\varepsilon t})$,


а последнее --- $O ((t+1)^ {\alpha-k-1})$. 


Тем самым доказано и


первое утверждение леммы.


\end{pf}








\begin{lem}


Для любого комплексного~$\alpha$, действительного~$\gamma\neq 0$ 


и натурального~$k$


$$


\int_0^t e^{i\gamma s}(s+1)^{\alpha}ds=c_k+e^{i\gamma t}(t+1)^{\alpha}


u (t),\ \ u \in A_k,\ \ c_k=\const.


$$


\end{lem}





{\bf Доказательство.} Пусть $k>\re \alpha$. Интегрируя по


частям $k+2$ раза, получим


$$


\int_0^t e^{i \gamma s}(s+1)^{\alpha}ds=e^{i \gamma t}


\sum\limits_{j=0}^{k+1} a_j (t+1)^{\alpha-j}+b_0+


b_1 \bigg[\int_0^\infty e^{i \gamma s}(s+1)^{\alpha-k-2}ds-


\int_t^\infty e^{i \gamma s}(s+1)^{\alpha-k-2}ds \bigg].


$$


Осталось заметить, что


$$


\bigg|\int_t^\infty e^{i \gamma s}(s+1)^{\alpha-k-2}ds \bigg|\le


\int_t^\infty (s+1)^{\re \alpha-k-2}ds=


c (t+1)^{\re \alpha-k-1}.\quad\square


$$





\begin{pf*}{Доказательство теоремы} Покажем, что для решения


$x (t)$ уравнения~(\ref{5}), эквивалентного~(\ref{1}), справедлива


оценка


\begin{equation}


\label{9}


x(t)=O (1)e^{\lambda t}(t+1)^{\alpha},\; \alpha=


\max\limits_{1\le j\le r} \re \alpha_j.


\end{equation}


Положим


$z (t)=e^{-\lambda t}(t+1)^{-\alpha} x(t)$.


Тогда


$z(t)=\int\limits_0^t Q(t,s) z(s)ds+g(t)$,


где $g(t)=O(1)$, $q_l (t)=e^{-\lambda t} h_l (t)=


o(e^{-\varepsilon t})$ и


\begin{gather*}


Q(t,s)=\sum\limits_{l=1}^n \bigg(\frac {s+1}{t+1} \bigg)^{\alpha}


q_l (t-s)\frac{\gamma_l (s)}{s+1}+\\


%


+\sum\limits_{j=1}^r \sum\limits_{l=1}^n e^{(\lambda_j-\lambda)t}


(t+1)^{\alpha_{j}-\alpha} (s+1)^{\alpha}


u_j (t) \frac{\gamma_l (s)}{s+1}


\int_{s}^t q_l (\tau-s)


e^{(\lambda-\lambda_j)\tau} (\tau+1)^{-\alpha_{j}}


\frac{v_j (\tau)}{\tau+1}d\tau.


\end{gather*}


Используя лемму~2, найдем


\begin{gather*}


\int_T^t |Q(t, s)|ds \le c \bigg[ (t+1)^{-\alpha}


\int_0^t e^{-\varepsilon(t-s)} (s+1)^{\alpha-1}ds+\\


%


+\sum\limits_{j=1}^r (t+1)^{\re \alpha_{j}-\alpha}


\int_T^t \int_0^\tau (s+1)^{\alpha-1}


e^{-\varepsilon(\tau-s)} (\tau+1)^{-{\rm Re}\alpha_{j}-1} 


ds d\tau \bigg] \le


\frac {c_1}{t+1}+\frac {c_1}{T+1}\le \frac {2c_1}{T+1}<1


\end{gather*}


при достаточно большом $T$. Следовательно (см., напр., [8]), ядро $Q$ 


устойчиво, и потому


$z (t)=O (1)$, что и доказывает~(\ref{9}).





Покажем теперь, что при любом целом $p\ge0$


\begin{equation}


\label{10}


x(t)=\sum\limits_{j=1}^{r} e^{\lambda_jt}(t+1)^{\alpha_{j}} \bigg[


\sum\limits_{k=0}^{p-1}\frac{x_{jk}}{(t+1)^{k}}+


O \bigg(\frac1{(t+1)^{p-1+\delta}} \bigg) \bigg],


\end{equation}


где $\delta\in(0,1]$ и таково, что


$\re \alpha_{j}-\re \alpha_{l}+\delta$ при всех $j$, $l$ не является


целым числом. При $p=0$ равенство~(\ref{10}) справедливо, т.\,к. если $j_0$ 


таково, что $\re \alpha_{j_0}=\alpha$, то согласно~(\ref{9})


$$


x (t)=e^{\lambda_{j_0}t}(t+1)^{\alpha_{j_0}}O (1)+


\sum\limits_{j\neq j_0}e^{\lambda_jt}(t+1)^{\alpha_{j}}\cdot 0.


$$





Пусть~(\ref{10}) справедливо при некотором $p\ge0$. Очевидно,


$h_l (t)=e^{\lambda_jt}e^{-\varepsilon t}o (1)$. 


Поэтому в силу леммы~2


\begin{equation}


\label{11}


\int_0^t H (t,s)x(s)ds=\sum\limits_{j=1}^{r}


e^{\lambda_jt}(t+1)^{\alpha_{j}} \bigg[


\sum\limits_{k=1}^{p}\frac{y_{jk}}{(t+1)^{k}}+


O \bigg(\frac1{(t+1)^{p+\delta}} \bigg) \bigg].


\end{equation}


Далее


\begin{gather*}


g_{jml} (t)=\int_0^t G_j (t,s)


e^{\lambda_m s}(s+1)^{\alpha_{m}-l}ds=\\


%


=e^{\lambda_jt}(t+1)^{\alpha_{j}}


u_j (t) \int_0^t e^{ (\lambda_m-\lambda_j)s}


(s+1)^{\alpha_{m}-\alpha_{j}-l-1}v_j (s)ds.


\end{gather*}





Пусть $k$ --- достаточно большое натуральное число. Тогда при $m\neq j$


по лемме~3


\begin{equation}


\label{12}


g_{jml} (t)=c_{jml} e^{\lambda_jt}(t+1)^{\alpha_{j}}


u_j (t)+e^{\lambda_m t}(t+1)^{\alpha_{m}-l-1}


v_{jml} (t),\ \ v_{jml} \in A_{k},


\end{equation}


а


\begin{equation}


\label{13}


g_{jjl} (t)=e^{\lambda_j t}(t+1)^{\alpha_{j}}


v_{jjl} (t),\ \ v_{jjl} \in A_{k}.


\end{equation}


Наконец,


\begin{multline}


\label{14}


\int_0^t G_j(t,s) e^{\lambda_m s} (s+1)^{\alpha_{m}-p-


\delta}O (1)ds=\\


%


=e^{\lambda_jt}(t+1)^{\alpha_{j}} u_j (t)


\int_0^t (s+1)^{\alpha_{m}-\alpha_{j}-p-\delta-1}O (1)ds=\\


%


=c_{jm}e^{\lambda_jt}(t+1)^{\alpha_{j}} u_j (t)


+e^{\lambda_m t}(t+1)^{\alpha_{m}-p-\delta}O (1),


\end{multline}


т.\,к. если $\re \alpha_{m}-\re \alpha_{j}-p-\delta<0$, то


$$


\int_0^t (s+1)^{\alpha_{m}-\alpha_{j}-p-\delta-1}O (1)ds=


c_{jm}+\int_t^\infty (s+1)^{\alpha_{m}-\alpha_{j}-p-\delta-1}


O (1)ds=


c_{jm}+(t+1)^{\alpha_{m}-\alpha_{j}-p-\delta} O (1),


$$


а если $\re \alpha_{m}-\re \alpha_{j}-p-\delta>0$, то


$$


\int_0^t (s+1)^{\alpha_{m}-\alpha_{j}-p-\delta-1}O (1)ds


=(t+1)^{\alpha_{m}-\alpha_{j}-p-\delta} O (1).


$$


Из~(\ref{5}), (\ref{10})--(\ref{14}) следует


\begin{equation}


\label{15}


x(t)=\sum\limits_{j=1}^{r} e^{\lambda_jt}(t+1)^{\alpha_{j}} \bigg[


\sum\limits_{k=0}^{p}\frac{c_{jk}}{(t+1)^{k}}+


O \bigg(\frac1{(t+1)^{p+\delta}} \bigg) \bigg].


\end{equation}


По индукции заключаем, что~(\ref{10}) справедливо при всех $p$.





При переходе от~(\ref{10}) к~(\ref{15}) слагаемое 


$\frac {x_{jk}} { (t+1)^{k-\alpha_j}}$


дает вклад в различные слагаемые разложения~(\ref{15}),


поэтому не ясно, совпадают ли $c_{jk}$ с соответствующими $x_{jk}$. 


Покажем, что с ростом $p$ коэффиценты разложения все же стабилизируются, 


т.\,е. каждый коэффициент $c_{jk}=x_{jk}$ при достаточно больших $p$. 


Допустим противное. Вычитая из~(\ref{15}) равенство~(\ref{10}), 


найдем (при достаточно большом $p$)


\begin{equation}


\label{16}


\sum\limits_{j=1}^r e^{\lambda_jt}(t+1)^{\alpha_{j}}z_j (t)\equiv 0,


\end{equation}


где ${z_j\in A_p}$ и присутствуют $z_j$ с ненулевыми коэффициентами 


разложения. Пусть $l_j$ --- номер первого ненулевого коэффициента разложения 


$z_j (t)$, если такой есть, и $l_j=\infty$ в противном случае. Меняя, 


если необходимо, нумерацию, можем считать, что 


$\re (\alpha_1-l_1)\ge \re (\alpha_2-l_2)\ge 


\dotsb \ge\re (\alpha_r-l_r)$. Разделив~(\ref{16}) на 


$e^{\lambda t}(t+1)^{\re (\alpha_{1}-l_1)}$, получим


$$


\sum\limits_{j=1}^k


e^{i\gamma_jt}(t+1)^{i\beta_{j}} (c_j+o (1))=o (1),


$$


где все $c_j\neq 0$ и $\gamma_j=\mathrm{Im}\,\lambda_j$ попарно различны. 


Отсюда следует


$$


\sum\limits_{j=1}^k e^{i\gamma_jt}(t+1)^{i\beta_{j}}c_j=o (1).


$$


Покажем по индукции, что это возможно лишь при $c_j=0$. Действительно, если 


$k=1$, то это очевидно. Пусть это верно при всех $k\le m$. Заменяя в равенстве


\begin{equation}


\label{17}


\sum\limits_{j=1}^{m+1}


e^{i\gamma_jt}(t+1)^{i\beta_{j}}c_j=o (1)


\end{equation}


$t$ на $t+\varepsilon$ и пользуясь тем, что 


$(t+\varepsilon+1)^{i\beta_{j}}\sim(t+1)^{i\beta_{j}}$ 


при $t \to \infty$, получим


$$


\sum\limits_{j=1}^{m+1}


e^{i\gamma_jt}(t+1)^{i\beta_{j}}e^{i\gamma_j \varepsilon }c_j=o (1).


$$


Разделим это равенство на 


$e^{i\gamma_{m+1}\varepsilon}$


и вычтем результат из~(\ref{17}). Тогда


$$


\sum\limits_{j=1}^{m}


e^{i\gamma_jt}(t+1)^{i\beta_{j}}c_j


(1-e^{i\varepsilon (\gamma_j-\gamma_{m+1})})=o (1).


$$


Отсюда по предположению индукции $c_j=0$, $j=1,\dotsc,m$, т.\,к.


существует $\varepsilon$, при котором 


$ e^{i\varepsilon (\gamma_j-\gamma_{m+1})}\neq1$, $j=1,\dotsc,m$.


Но тогда из~(\ref{17}) следует, что и $c_{m+1}=0$. Таким образом, 


равенство~(\ref{17}) возможно лишь при $c_j=0$, что и доказывает 


стабилизируемость коэффициентов разложения в~(\ref{10}) с ростом $p$.


\end{pf*}





\begin{thebibliography}{9}





\bibitem{1}


Дербенев В.А., Цалюк З.Б. 


{\it Асимптотика резольвенты неустойчивого уравнения Вольтерра 


с разностным ядром}


// Матем. заметки. -- 1997. -- Т.\,62. -- Вып.\,1. -- C.\,88--94.


\bibitem{2}


Цалюк З.Б.


{ \it Асимптотическая структура резольвенты неустойчивого


уравнения Вольтерра с разностным ядром}


// Изв. вузов. Математика. -- 2000. -- \No\,4. -- C.\,50--55.


\bibitem{3}


Цалюк З.Б.


{\it Структура резольвенты системы уравнений с


разностным ядром}


// Изв. вузов. Математика. -- 2001. -- \No\,6. -- C.\,71--80.


\bibitem{4}


Цалюк З.Б. 


{\it Структура резольвенты уравнения восстановления}


// Изв. вузов. Сев.-Кавк. регион. Естеств. науки. -- 2001.


-- C.\,150--151.


\bibitem{5}


Дербенев В.А., Пуляев В.Ф. 


{\it Структура резольвенты и устойчивость линейных интегральных и 


интегродифференциальных уравнений}


// Изв. вузов. Сев.--Кавк. научн. центр высшей школы.


Естеств. науки. -- 1992. -- \No\,1--2. -- C.\,7--14.


\bibitem{6}


Пуляев В.Ф., Цалюк З.Б. 


{\it Об асимптотическом поведении решений интегральных уравнений 


Вольтерра в банаховых пространствах}


// Изв. вузов. Математика. -- 1991. -- \No\,12. -- C.\,47--55.


\bibitem{7}


Коддингтон Э.А., Левинсон Н. 


{\it Теория обыкновенных дифференциальных уравнений}.


-- М: Ин. лит., 1958. -- 474~с.


\bibitem{8}


Цалюк З.Б.


{\it Интегральные уравнения Вольтерра}


// Итоги науки и техн. Матем. анализ. -- М.: ВИНИТИ, 1977. -- Т. 15. -- 


C.\,131--198.





\end{thebibliography}





\vskip 2cm





\post{\it Кубанский государственный}{\it Поступила}


\post{\it университет $($г.\,Краснодар$)$}{24.01.2002}





\label{end}


\end{document}















































